\documentclass[11pt]{article}
\usepackage[a4paper,margin=1in]{geometry}
\usepackage{amsmath,amssymb,amsthm}
\usepackage[T1]{fontenc}
\usepackage{lmodern}
\usepackage[hidelinks]{hyperref}
\setlength{\emergencystretch}{3em}

\theoremstyle{plain}
\newtheorem{theorem}{Theorem}
\newtheorem{lemma}[theorem]{Lemma}
\theoremstyle{definition}
\newtheorem{definition}[theorem]{Definition}
\theoremstyle{remark}
\newtheorem{remark}[theorem]{Remark}

\DeclareMathOperator{\rank}{rank}
\DeclareMathOperator{\rankp}{rank_+}

\title{A Counterexample to Conjecture 00000002850:\\
Nonnegative Rank Separates from Rank Already at Rank $3$}
\author{%
  Feiyu\thanks{Independent researcher.}
  \and
  Claude (Opus 5.5)\thanks{Anthropic.}%
}
\date{2026-10-02}

\begin{document}
\maketitle

\begin{abstract}
Conjecture 00000002850 asserts that the separation of nonnegative rank from
ordinary rank ``starts at rank $4$''. It starts at rank $3$: the $4\times4$
slack matrix of a square has rank $3$ and nonnegative rank $4$. The lower bound
on the nonnegative rank is a four-element fooling set. Everything is formalized
in Lean~4 against Mathlib.
\end{abstract}

\section{The conjecture as stated}

The original text of conjecture 00000002850 reads:

\begin{quote}
Definition: Nonnegative matrix factorization (NMF): the algebraic geometry of
nonnegativity constraints. Conjecture: Rank rigidity of NMF: the separation of
nonnegative rank from ordinary rank starts at rank 4, with the minimal
separating matrix an explicit $4\times4$ nonnegative matrix; the geometric
characterization of the separation is a non-triangulable pair of nested
polytopes.
\end{quote}

\section{Reading of the statement}

\begin{definition}
The \emph{nonnegative rank} $\rankp(M)$ of an entrywise nonnegative real
$m\times n$ matrix $M$ is the least $k$ such that $M=WH$ with $W$ ($m\times k$)
and $H$ ($k\times n$) entrywise nonnegative. Always $\rankp(M)\ge\rank(M)$. We
say the two ranks \emph{separate at rank $r$} if some nonnegative matrix of
rank $r$ has nonnegative rank larger than $r$.
\end{definition}

``The separation starts at rank $4$'' asserts that $4$ is the least $r$ at which
the ranks separate: they separate at rank $4$, and at no $r<4$. The statement
distinguishes this rank from the size of the minimal example (``an explicit
$4\times4$ matrix''), so the number $4$ in ``starts at rank $4$'' refers to the
rank. We refute it by exhibiting a separation at rank $3$.

\section{Result}

\begin{theorem}\label{thm:main}
The slack matrix of a square
\[
  S=\begin{pmatrix}1&1&0&0\\0&1&1&0\\0&0&1&1\\1&0&0&1\end{pmatrix}
\]
has $\rank(S)=3$ and $\rankp(S)=4$. Hence the ranks separate at rank $3$, and
Conjecture 00000002850 is false.
\end{theorem}

\section{Proof}

\paragraph{The rank is $3$.} The fourth row is $r_1-r_2+r_3$, so $\rank S\le3$,
and the leading $3\times3$ block
$\left(\begin{smallmatrix}1&1&0\\0&1&1\\0&0&1\end{smallmatrix}\right)$ has
determinant $1$, so $\rank S\ge3$.

\paragraph{The nonnegative rank is $4$.} Trivially $\rankp S\le4$. Suppose
$S=WH$ with $W,H\ge0$ and inner dimension $3$, so $S_{ij}=\sum_{k=1}^3W_{ik}H_{kj}$
is a sum of nonnegative terms. Each diagonal entry $S_{ii}=1$ is positive, so
for each $i$ there is a $k=f(i)$ with $W_{ik}>0$ and $H_{ki}>0$. If $f(i)=f(j)=k$
for some $i\ne j$, then $S_{ij}\ge W_{ik}H_{kj}>0$ and $S_{ji}\ge W_{jk}H_{ki}>0$.
But for every pair $i\ne j$ at least one of $S_{ij},S_{ji}$ is $0$; the diagonal
positions form a fooling set. So $f\colon\{1,2,3,4\}\to\{1,2,3\}$ is injective,
which is impossible.

\section{Formalization}

The accompanying Lean~4 project (\texttt{lean/Submission/Basic.lean}) states
the conjecture and proves its negation against Mathlib.

\begin{itemize}
  \item \texttt{HasNonnegFactorization M k} --- $M=WH$ with $W,H$ entrywise
    nonnegative and inner dimension $k$.
  \item \texttt{SeparatesAt r} --- some entrywise nonnegative real matrix $M$
    has \texttt{M.rank = r} (Mathlib's rank) and no nonnegative factorization
    of inner dimension $r$.
  \item \texttt{ConjectureHolds} --- \texttt{SeparatesAt 4} and
    $\lnot\,$\texttt{SeparatesAt r} for all $r<4$.
  \item \texttt{S}, \texttt{rank\_S} --- the matrix and its rank, via an
    explicit factorization through dimension $3$ and the invertible leading
    block.
  \item \texttt{S\_fooling}, \texttt{not\_hasNonnegFactorization\_S} --- the
    fooling-set argument.
  \item \texttt{separatesAt\_three}, \texttt{conjecture\_00000002850\_false}.
\end{itemize}

The toolchain is \texttt{leanprover/lean4:v4.33.1}, Mathlib is pinned to
\texttt{v4.33.1}, and \texttt{lake build} succeeds. The project contains no
\texttt{sorry}, no \texttt{admit} and no \texttt{native\_decide}, and
introduces no axioms:
\begin{center}
\texttt{\#print axioms conjecture\_00000002850\_false}
\end{center}
reports exactly \texttt{propext}, \texttt{Classical.choice} and
\texttt{Quot.sound}.

\section{Remarks}

For matrices of rank at most $2$ the two ranks coincide \cite{cr}, so the
separation starts exactly at rank $3$, with $S$ as a minimal example. Geometrically,
$S$ is the slack matrix of a square, and $\rankp(S)=4$ reflects the fact that a
square in the plane cannot be written as the projection of a polytope with
fewer than $4$ facets, i.e.\ its extension complexity is $4$ \cite{yannakakis}.
The ``nested polytopes'' picture of the statement belongs to rank $3$, not $4$:
a rank-$3$ matrix has nonnegative rank $3$ exactly when a triangle can be
nested between the corresponding pair of polygons.

\begin{thebibliography}{9}
\bibitem{cr} J.~E.~Cohen and U.~G.~Rothblum, \emph{Nonnegative ranks,
  decompositions, and factorizations of nonnegative matrices}, Linear Algebra
  Appl. 190 (1993), 149--168.
\bibitem{yannakakis} M.~Yannakakis, \emph{Expressing combinatorial
  optimization problems by linear programs}, J. Comput. System Sci. 43 (1991),
  441--466.
\end{thebibliography}

\end{document}
